%==========================================================================
% Section: Poisson-gamma / binomial-probit thinning model
%          for bivariate claim frequency
%==========================================================================
\section{Binomial thinning model for bivariate claim frequency}
\label{sec.thin}

In this section, we develop a model for two nested claim counts.
For each region \(i\) and period \(t\), let \(N_{it}\) denote the total
number of claims and let \(M_{it}\) denote the number of extreme claims
among them, for example, the claims whose amounts exceed a given
threshold. By construction,
\[
0\le M_{it}\le N_{it}.
\]
A natural way to respect this constraint is to regard \(M_{it}\) as a
binomial thinning of \(N_{it}\): each of the \(N_{it}\) claims is
classified as extreme or not. Repeated observations of the same region
over time are linked through two region-specific random effects, one
for the claim frequency and one for the probability that a claim is
extreme. Throughout this section, the two random effects are assumed to
be independent and static over time, and predictions are obtained as
exact Bayesian posterior means.

\begin{model}[Poisson-gamma binomial-probit thinning model]
\label{model.thin}
Let
\[
r,\sigma\in\Real_{>0},
\qquad
\alpha,\beta\in\Real^p,
\]
and let \(x_{it}\in\Real^p\), \(i=1,\ldots,I\), \(t\in\mathbb N\), be
predetermined exogenous covariates, which are allowed to vary over time.
Define
\[
\lambda_{it}:=\exp\left(x_{it}^\top\alpha\right)\in\Real_{>0},
\qquad
\eta_{it}:=x_{it}^\top\beta\in\Real.
\]
We define the bivariate sequences
\(\left((N_{it},M_{it})\right)_{t\in\mathbb N}\), \(i=1,\ldots,I\),
through the following random effect assumption and observation
equations.
\begin{itemize}
  \item[i.] Random effects: For each region \(i\), assume
  \[
  \Theta_i^{(1)}\sim\operatorname{Gamma}(r,r),
  \qquad
  \Theta_i^{(2)}\sim\operatorname{N}(0,\sigma^2),
  \]
  where \(\Theta_i^{(1)}\) and \(\Theta_i^{(2)}\) are independent.
  Furthermore, the pairs
  \(\left(\Theta_i^{(1)},\Theta_i^{(2)}\right)\), \(i=1,\ldots,I\),
  are mutually independent.

  \item[ii.] Observation equations: Conditional on
  \(\left(\Theta_i^{(1)},\Theta_i^{(2)}\right)\), assume that the pairs
  \((N_{it},M_{it})\), \(t\in\mathbb N\), are mutually independent with
  \[
  N_{it}\mid\Theta_i^{(1)},\Theta_i^{(2)}
  \sim
  \operatorname{Pois}\left(\lambda_{it}\Theta_i^{(1)}\right)
  \]
  and
  \[
  M_{it}\mid N_{it},\Theta_i^{(1)},\Theta_i^{(2)}
  \sim
  \operatorname{Bin}\left(N_{it},\Phi\left(\eta_{it}+\Theta_i^{(2)}\right)\right),
  \]
  Furthermore, the regions are independent: the sequences
  \(\left((N_{it},M_{it})\right)_{t\in\mathbb N}\), \(i=1,\ldots,I\),
  are mutually independent.
\end{itemize}
\end{model}
Here, \(\operatorname{Bin}(0,p)\) is the point mass at \(0\), so that
\(M_{it}=0\) whenever \(N_{it}=0\).

Since \(\E{\Theta_i^{(1)}}=1\) and \(\E{\Theta_i^{(2)}}=0\), the
intercepts in \(\alpha\) and \(\beta\) are identifiable, and
\(\lambda_{it}\) is the a priori expected total claim count. The random
effect \(\Theta_i^{(1)}\) captures the unobserved heterogeneity of the
claim frequency across regions, whereas \(\Theta_i^{(2)}\) captures the
unobserved heterogeneity of the propensity of a claim to be extreme.
For notational convenience, write
\[
p_{it}(\theta):=\Phi\left(\eta_{it}+\theta\right),
\qquad \theta\in\Real.
\]

\begin{remark}\label{rem.thin.exposure}
All results in this section use \(\lambda_{it}\) and \(\eta_{it}\) only
as known constants. Hence they remain valid for any
\(\lambda_{it}\in\Real_{>0}\) and \(\eta_{it}\in\Real\); in particular,
an exposure offset \(e_{it}>0\) can be included through
\(\lambda_{it}=e_{it}\exp\left(x_{it}^\top\alpha\right)\).
\end{remark}

The following lemma is the well-known thinning property of the Poisson
distribution. It shows that, conditional on the random effects, the
extreme and non-extreme claim counts are independent Poisson random
variables.

\begin{lemma}[Poisson thinning]
\label{lem.thin}
Under Model~\ref{model.thin}, conditional on
\(\left(\Theta_i^{(1)},\Theta_i^{(2)}\right)\), the random variables
\(M_{it}\) and \(N_{it}-M_{it}\) are independent with
\[
M_{it}\mid\Theta_i^{(1)},\Theta_i^{(2)}
\sim
\operatorname{Pois}\left(
\lambda_{it}\Theta_i^{(1)}p_{it}\left(\Theta_i^{(2)}\right)
\right)
\]
and
\[
N_{it}-M_{it}\mid\Theta_i^{(1)},\Theta_i^{(2)}
\sim
\operatorname{Pois}\left(
\lambda_{it}\Theta_i^{(1)}
\left\{1-p_{it}\left(\Theta_i^{(2)}\right)\right\}
\right).
\]
\end{lemma}

\begin{proof}
Fix \(\theta_1>0\) and \(\theta_2\in\Real\), and write
\(\Lambda:=\lambda_{it}\theta_1\) and \(p:=p_{it}(\theta_2)\). For
\(m,k\in\mathbb N_0\), Model~\ref{model.thin} gives
\[
\begin{aligned}
&\P{M_{it}=m,\,N_{it}-M_{it}=k\mid
\Theta_i^{(1)}=\theta_1,\Theta_i^{(2)}=\theta_2}\\
&\quad=
\P{N_{it}=m+k\mid\Theta_i^{(1)}=\theta_1}
\binom{m+k}{m}p^m(1-p)^k\\
&\quad=
\frac{e^{-\Lambda}\Lambda^{m+k}}{(m+k)!}
\frac{(m+k)!}{m!\,k!}p^m(1-p)^k\\
&\quad=
\frac{e^{-\Lambda p}(\Lambda p)^m}{m!}
\cdot
\frac{e^{-\Lambda(1-p)}\left\{\Lambda(1-p)\right\}^k}{k!},
\end{aligned}
\]
which proves the result.
\end{proof}

Since the regions are independent, it suffices to study a single
region. The following remark collects the notation for its observed
history.

\begin{remark}\label{rem.thin.notation}
Fix a region \(i\) and \(t\in\mathbb N\), and consider a realized
history
\[
\mathcal H_{it}:=
\left\{
(N_{ij},M_{ij})=(n_{ij},m_{ij}):\ j=1,\ldots,t
\right\},
\qquad
n_{ij},m_{ij}\in\mathbb N_0,
\quad
m_{ij}\le n_{ij}.
\]
Define the cumulative claim count and the cumulative a priori
frequency by
\[
n_{i\bullet t}:=\sum_{j=1}^t n_{ij},
\qquad
\lambda_{i\bullet t}:=\sum_{j=1}^t \lambda_{ij}.
\]

Each of the \(n_{i\bullet t}\) observed claims is regarded as a
Bernoulli trial, which is a success if the claim is extreme. To this
end, for \(j=1,\ldots,t\), define the sign vector
\[
u^{(j)}:=
\big(
\underbrace{1,\ldots,1}_{m_{ij}},
\underbrace{-1,\ldots,-1}_{n_{ij}-m_{ij}}
\big)^\top
\in\{-1,1\}^{n_{ij}},
\]
which is empty when \(n_{ij}=0\), and stack these vectors as
\[
u_{1:n_{i\bullet t}}:=
\left(
u^{(1)\top},\ldots,u^{(t)\top}
\right)^\top,
\qquad
h_{1:n_{i\bullet t}}:=
\big(
\underbrace{\eta_{i1},\ldots,\eta_{i1}}_{n_{i1}},
\ldots,
\underbrace{\eta_{it},\ldots,\eta_{it}}_{n_{it}}
\big)^\top .
\]
That is, \(u_k\) is the sign of the \(k\)th claim and \(h_k\) is the
value of \(\eta_{ij}\) in the period \(j\) of the \(k\)th claim. We then
define
\[
\mu_{it}:=
\left(u_1h_1,\ldots,u_{n_{i\bullet t}}h_{n_{i\bullet t}}\right)^\top
\in\Real^{n_{i\bullet t}},
\qquad
\boldsymbol U_{it}:=
u_{1:n_{i\bullet t}}u_{1:n_{i\bullet t}}^\top
\in\Real^{n_{i\bullet t}\times n_{i\bullet t}},
\]
which are the mean vector and the rank-one matrix appearing in
Lemma~\ref{lem.bp} for the trials formed by the observed claims.

For \(n,m\in\mathbb N_0\) with \(m\le n\), let
\(\mu_{it}^{(n,m)}\in\Real^{n_{i\bullet t}+n}\) and
\(\boldsymbol U_{it}^{(n,m)}\in\Real^{(n_{i\bullet t}+n)\times(n_{i\bullet t}+n)}\)
denote the corresponding quantities for the augmented history
\(\mathcal H_{it}\cup\{(N_{i,t+1},M_{i,t+1})=(n,m)\}\). That is,
\(n\) additional trials with \(\eta\)-value \(\eta_{i,t+1}\) are appended,
\(m\) of them with sign \(1\) and \(n-m\) of them with sign \(-1\).

Finally, define
\[
D_{it}(\theta):=
\prod_{j=1}^t
p_{ij}(\theta)^{m_{ij}}
\left\{1-p_{ij}(\theta)\right\}^{n_{ij}-m_{ij}},
\qquad\theta\in\Real.
\]
Since \(1-\Phi(a)=\Phi(-a)\) for \(a\in\Real\), we have
\begin{equation}\label{eq.thin.D}
D_{it}(\theta)
=
\prod_{k=1}^{n_{i\bullet t}}
\Phi\left(u_k(h_k+\theta)\right).
\end{equation}
For \(n_{i\bullet t}=0\), we adopt the conventions
\(D_{it}(\theta)=1\) and
\(\Phi_0\left(\mu_{it};\boldsymbol I_0+\sigma^2\boldsymbol U_{it}\right)=1\).
\end{remark}

The following theorem shows that the joint probability of the observed
history factorizes into a closed-form Poisson-gamma part, which depends
only on \((\alpha,r)\), and a Bernoulli-probit part of the form in
Lemma~\ref{lem.bp}, which depends only on \((\beta,\sigma)\).

\begin{theorem}[Likelihood factorization]
\label{thm.thin.lik}
Consider Model~\ref{model.thin} with the realized history
\(\mathcal H_{it}\), and let the notation be as in
Remark~\ref{rem.thin.notation}. Then
\begin{equation}\label{eq.thin.factorization}
\P{\mathcal H_{it}}
=
L^{(1)}_{it}\,L^{(2)}_{it},
\end{equation}
where
\begin{equation}\label{eq.thin.L1}
L^{(1)}_{it}
:=
\P{N_{i,1:t}=n_{i,1:t}}
=
\frac{\Gamma\left(r+n_{i\bullet t}\right)}
{\Gamma(r)\prod_{j=1}^t n_{ij}!}
\,
\frac{r^r\prod_{j=1}^t\lambda_{ij}^{n_{ij}}}
{\left(r+\lambda_{i\bullet t}\right)^{r+n_{i\bullet t}}}
\end{equation}
and
\begin{equation}\label{eq.thin.L2}
\begin{aligned}
L^{(2)}_{it}
&:=
\P{M_{i,1:t}=m_{i,1:t}\mid N_{i,1:t}=n_{i,1:t}}\\
&=
\prod_{j=1}^t\binom{n_{ij}}{m_{ij}}
\int_{\Real}
\frac{1}{\sigma}\phi\left(\frac{\theta}{\sigma}\right)
D_{it}(\theta)
\,{\rm d}\theta\\
&=
\prod_{j=1}^t\binom{n_{ij}}{m_{ij}}
\,
\Phi_{n_{i\bullet t}}
\left(
\mu_{it};
\boldsymbol I_{n_{i\bullet t}}
+
\sigma^2\boldsymbol U_{it}
\right).
\end{aligned}
\end{equation}
\end{theorem}

\begin{proof}
By the conditional independence in Model~\ref{model.thin},
\[
\P{\mathcal H_{it}}
=
\E{
\prod_{j=1}^t
\P{N_{ij}=n_{ij}\mid\Theta_i^{(1)}}
\P{M_{ij}=m_{ij}\mid N_{ij}=n_{ij},\Theta_i^{(2)}}
}.
\]
The first factor in the product depends only on \(\Theta_i^{(1)}\) and
the second only on \(\Theta_i^{(2)}\). Since \(\Theta_i^{(1)}\) and
\(\Theta_i^{(2)}\) are independent, we obtain
\[
\P{\mathcal H_{it}}
=
\E{
\prod_{j=1}^t
\P{N_{ij}=n_{ij}\mid\Theta_i^{(1)}}
}
\,
\E{
\prod_{j=1}^t
\P{M_{ij}=m_{ij}\mid N_{ij}=n_{ij},\Theta_i^{(2)}}
}.
\]

For the first expectation, the gamma integral gives
\[
\begin{aligned}
\E{
\prod_{j=1}^t
\P{N_{ij}=n_{ij}\mid\Theta_i^{(1)}}
}
&=
\int_0^\infty
\prod_{j=1}^t
\frac{e^{-\lambda_{ij}\theta}(\lambda_{ij}\theta)^{n_{ij}}}{n_{ij}!}
\,
\frac{r^r}{\Gamma(r)}\theta^{r-1}e^{-r\theta}
\,{\rm d}\theta\\
&=
\frac{r^r\prod_{j=1}^t\lambda_{ij}^{n_{ij}}}
{\Gamma(r)\prod_{j=1}^t n_{ij}!}
\int_0^\infty
\theta^{r+n_{i\bullet t}-1}
e^{-(r+\lambda_{i\bullet t})\theta}
\,{\rm d}\theta,
\end{aligned}
\]
which equals the right-hand side of \eqref{eq.thin.L1}. Summing the
first display over all \(m_{i,1:t}\) with \(0\le m_{ij}\le n_{ij}\)
shows that this expectation is \(\P{N_{i,1:t}=n_{i,1:t}}\), because
the second expectation sums to one.

For the second expectation, we have
\[
\E{
\prod_{j=1}^t
\P{M_{ij}=m_{ij}\mid N_{ij}=n_{ij},\Theta_i^{(2)}}
}
=
\prod_{j=1}^t\binom{n_{ij}}{m_{ij}}
\E{D_{it}\left(\Theta_i^{(2)}\right)},
\]
which proves the second line of \eqref{eq.thin.L2}. It then follows
from \eqref{eq.thin.factorization} and \eqref{eq.thin.L1} that this
expression is \(\P{M_{i,1:t}=m_{i,1:t}\mid N_{i,1:t}=n_{i,1:t}}\).
Finally, by \eqref{eq.thin.D},
\[
\E{D_{it}\left(\Theta_i^{(2)}\right)}
=
\E{
\prod_{k=1}^{n_{i\bullet t}}
\Phi\left(u_k\left(h_k+\Theta_i^{(2)}\right)\right)
},
\]
and applying Lemma~\ref{lem.bp} with \(k=n_{i\bullet t}\),
\(u_{1:k}\), and \(h_{1:k}\) proves the last line of
\eqref{eq.thin.L2}.
\end{proof}

\begin{remark}[Separable estimation]
\label{rem.thin.estimation}
By Theorem~\ref{thm.thin.lik} and the independence of the regions,
the log-likelihood of the observations
\(\mathcal H_{1T},\ldots,\mathcal H_{IT}\) decomposes as
\[
\ell(\alpha,r,\beta,\sigma)
=
\ell^{(1)}(\alpha,r)+\ell^{(2)}(\beta,\sigma),
\qquad
\ell^{(k)}:=\sum_{i=1}^I\log L^{(k)}_{iT},
\quad k=1,2.
\]
Hence the maximum likelihood estimators of \((\alpha,r)\) and
\((\beta,\sigma)\) are obtained from two separate optimizations, and
the Fisher information matrix is block-diagonal. Here,
\(\ell^{(1)}\) is the log-likelihood of a negative binomial
random effect panel model for \(N_{it}\), and \(\ell^{(2)}\) is the
log-likelihood of a binomial probit model for \(M_{it}\) given
\(N_{it}\) with a Gaussian random intercept. Periods with \(n_{it}=0\)
do not contribute to \(\ell^{(2)}\).
\end{remark}

\begin{remark}[Computation]
\label{rem.thin.computation}
The dimension \(n_{i\bullet t}\) of the multivariate normal
distribution function in \eqref{eq.thin.L2} equals the cumulative
number of claims of region \(i\), which can be large. However, since
\(\boldsymbol U_{it}\) is of rank one, the second line of
\eqref{eq.thin.L2} expresses the same quantity as the one-dimensional
integral
\[
\Phi_{n_{i\bullet t}}
\left(
\mu_{it};
\boldsymbol I_{n_{i\bullet t}}+\sigma^2\boldsymbol U_{it}
\right)
=
\int_{\Real}\phi(z)\,D_{it}(\sigma z)\,{\rm d}z,
\]
which depends on the data only through the \(t\) pairs
\((n_{ij},m_{ij})\) and can be evaluated accurately by Gauss--Hermite
quadrature.
\end{remark}

The following corollary shows that the two random effects remain
independent a posteriori, and that the frequency random effect has a
conjugate gamma posterior.

\begin{corollary}[Posterior distribution of the random effects]
\label{cor.thin.posterior}
Under the setting of Theorem~\ref{thm.thin.lik}, the following
results hold.
\begin{itemize}
  \item[i.] Conditional on \(\mathcal H_{it}\), the random effects
  \(\Theta_i^{(1)}\) and \(\Theta_i^{(2)}\) are independent.

  \item[ii.] We have
  \[
  \Theta_i^{(1)}\mid\mathcal H_{it}
  \sim
  \operatorname{Gamma}\left(r+n_{i\bullet t},\,r+\lambda_{i\bullet t}\right).
  \]

  \item[iii.] The conditional distribution
  \(\Theta_i^{(2)}\mid\mathcal H_{it}\) has density
  \(\pi(\theta\mid\mathcal H_{it})\) at \(\theta\in\Real\) given by
  \begin{equation}\label{eq.thin.posterior2}
  \pi(\theta\mid\mathcal H_{it})
  =
  \frac{
  \displaystyle
  \frac{1}{\sigma}\phi\left(\frac{\theta}{\sigma}\right)
  \prod_{k=1}^{n_{i\bullet t}}
  \Phi\left(u_k(h_k+\theta)\right)
  }{
  \Phi_{n_{i\bullet t}}
  \left(
  \mu_{it};
  \boldsymbol I_{n_{i\bullet t}}+\sigma^2\boldsymbol U_{it}
  \right)
  }.
  \end{equation}
\end{itemize}
\end{corollary}

\begin{proof}
Let \(\pi_1\) and \(\pi_2\) denote the prior densities of
\(\Theta_i^{(1)}\) and \(\Theta_i^{(2)}\), respectively. By Bayes'
formula and the proof of Theorem~\ref{thm.thin.lik}, the joint
posterior density of
\(\left(\Theta_i^{(1)},\Theta_i^{(2)}\right)\) at
\((\theta_1,\theta_2)\in\Real_{>0}\times\Real\) is
\[
\frac{
\pi_1(\theta_1)
\prod_{j=1}^t\P{N_{ij}=n_{ij}\mid\Theta_i^{(1)}=\theta_1}
}{L^{(1)}_{it}}
\cdot
\frac{
\pi_2(\theta_2)
\prod_{j=1}^t\binom{n_{ij}}{m_{ij}}D_{it}(\theta_2)
}{L^{(2)}_{it}}.
\]
Each factor integrates to one, which proves part i. The first factor
is proportional to
\[
\theta_1^{r+n_{i\bullet t}-1}e^{-(r+\lambda_{i\bullet t})\theta_1},
\]
which proves part ii. Substituting \eqref{eq.thin.D} and
\eqref{eq.thin.L2} into the second factor proves part iii.
\end{proof}

By part ii of Corollary~\ref{cor.thin.posterior}, the posterior of the
frequency random effect depends on the history only through
\(n_{i\bullet t}\); in particular, the extreme claim counts
\(m_{i,1:t}\) carry no information on \(\Theta_i^{(1)}\). Conversely,
by part iii, the total claim counts affect the posterior of
\(\Theta_i^{(2)}\) only by determining the number of Bernoulli trials.
The density \eqref{eq.thin.posterior2} is a Gaussian density
multiplied by a product of probit functions, normalized by the
multivariate normal distribution function in Lemma~\ref{lem.bp}.

We now present the exact Bayesian one-step-ahead predictions.

\begin{theorem}[Closed-form one-step prediction]
\label{thm.thin.pred}
Consider Model~\ref{model.thin} with the realized history
\(\mathcal H_{it}\), and let the notation be as in
Remark~\ref{rem.thin.notation}. Then, we have the following results.
\begin{itemize}
  \item[i.] For \(n,m\in\mathbb N_0\) with \(m\le n\), the predictive
  probability is
  \[
  \begin{aligned}
  &\P{N_{i,t+1}=n,\,M_{i,t+1}=m\mid\mathcal H_{it}}\\
  &\quad=
  \frac{\Gamma\left(r+n_{i\bullet t}+n\right)}
  {\Gamma\left(r+n_{i\bullet t}\right)n!}
  \left(
  \frac{r+\lambda_{i\bullet t}}{r+\lambda_{i\bullet t}+\lambda_{i,t+1}}
  \right)^{r+n_{i\bullet t}}
  \left(
  \frac{\lambda_{i,t+1}}{r+\lambda_{i\bullet t}+\lambda_{i,t+1}}
  \right)^{n}\\
  &\qquad\times
  \binom{n}{m}
  \frac{
  \Phi_{n_{i\bullet t}+n}
  \left(
  \mu_{it}^{(n,m)};
  \boldsymbol I_{n_{i\bullet t}+n}+\sigma^2\boldsymbol U_{it}^{(n,m)}
  \right)
  }{
  \Phi_{n_{i\bullet t}}
  \left(
  \mu_{it};
  \boldsymbol I_{n_{i\bullet t}}+\sigma^2\boldsymbol U_{it}
  \right)
  }.
  \end{aligned}
  \]

  \item[ii.] The one-step-ahead prediction of the total claim count is
  \[
  \E{N_{i,t+1}\mid\mathcal H_{it}}
  =
  \lambda_{i,t+1}\,
  \frac{r+n_{i\bullet t}}{r+\lambda_{i\bullet t}}
  =
  \lambda_{i,t+1}
  \left\{
  \zeta_{it}\frac{n_{i\bullet t}}{\lambda_{i\bullet t}}
  +
  \left(1-\zeta_{it}\right)
  \right\},
  \qquad
  \zeta_{it}:=\frac{\lambda_{i\bullet t}}{r+\lambda_{i\bullet t}}.
  \]

  \item[iii.] The predictive probability that a claim in period
  \(t+1\) is extreme is
  \[
  \widehat p_{i,t+1}
  :=
  \E{p_{i,t+1}\left(\Theta_i^{(2)}\right)\mid\mathcal H_{it}}
  =
  \frac{
  \Phi_{n_{i\bullet t}+1}
  \left(
  \mu_{it}^{(1,1)};
  \boldsymbol I_{n_{i\bullet t}+1}+\sigma^2\boldsymbol U_{it}^{(1,1)}
  \right)
  }{
  \Phi_{n_{i\bullet t}}
  \left(
  \mu_{it};
  \boldsymbol I_{n_{i\bullet t}}+\sigma^2\boldsymbol U_{it}
  \right)
  }.
  \]

  \item[iv.] For \(n\in\mathbb N_0\), the one-step-ahead prediction of
  the extreme claim count given the total claim count is
  \[
  \E{M_{i,t+1}\mid\mathcal H_{it},N_{i,t+1}=n}
  =
  n\,\widehat p_{i,t+1}.
  \]

  \item[v.] The one-step-ahead prediction of the extreme claim count is
  \[
  \E{M_{i,t+1}\mid\mathcal H_{it}}
  =
  \lambda_{i,t+1}\,
  \frac{r+n_{i\bullet t}}{r+\lambda_{i\bullet t}}\,
  \widehat p_{i,t+1}.
  \]
\end{itemize}
\end{theorem}

\begin{proof}
We first show part i. Applying Theorem~\ref{thm.thin.lik} to the
augmented history
\(\mathcal H_{it}\cup\{(N_{i,t+1},M_{i,t+1})=(n,m)\}\) and to
\(\mathcal H_{it}\), we obtain
\[
\P{N_{i,t+1}=n,\,M_{i,t+1}=m\mid\mathcal H_{it}}
=
\frac{L^{(1)}_{i,t+1}}{L^{(1)}_{it}}
\cdot
\frac{L^{(2)}_{i,t+1}}{L^{(2)}_{it}},
\]
where \(L^{(1)}_{i,t+1}\) and \(L^{(2)}_{i,t+1}\) are evaluated at the
augmented history. By \eqref{eq.thin.L1},
\[
\frac{L^{(1)}_{i,t+1}}{L^{(1)}_{it}}
=
\frac{\Gamma\left(r+n_{i\bullet t}+n\right)}
{\Gamma\left(r+n_{i\bullet t}\right)n!}
\,
\frac{\lambda_{i,t+1}^{n}\left(r+\lambda_{i\bullet t}\right)^{r+n_{i\bullet t}}}
{\left(r+\lambda_{i\bullet t}+\lambda_{i,t+1}\right)^{r+n_{i\bullet t}+n}},
\]
and, by the last line of \eqref{eq.thin.L2},
\[
\frac{L^{(2)}_{i,t+1}}{L^{(2)}_{it}}
=
\binom{n}{m}
\frac{
\Phi_{n_{i\bullet t}+n}
\left(
\mu_{it}^{(n,m)};
\boldsymbol I_{n_{i\bullet t}+n}+\sigma^2\boldsymbol U_{it}^{(n,m)}
\right)
}{
\Phi_{n_{i\bullet t}}
\left(
\mu_{it};
\boldsymbol I_{n_{i\bullet t}}+\sigma^2\boldsymbol U_{it}
\right)
},
\]
which proves part i.

To show part ii, note that, conditional on
\(\left(\Theta_i^{(1)},\Theta_i^{(2)}\right)\), \(N_{i,t+1}\) is
independent of \(\mathcal H_{it}\). Hence, by part ii of
Corollary~\ref{cor.thin.posterior},
\[
\E{N_{i,t+1}\mid\mathcal H_{it}}
=
\lambda_{i,t+1}\E{\Theta_i^{(1)}\mid\mathcal H_{it}}
=
\lambda_{i,t+1}\frac{r+n_{i\bullet t}}{r+\lambda_{i\bullet t}},
\]
and the second expression follows from
\(
(r+n_{i\bullet t})/(r+\lambda_{i\bullet t})
=
\zeta_{it}\,n_{i\bullet t}/\lambda_{i\bullet t}+1-\zeta_{it}.
\)

To show part iii, part iii of Corollary~\ref{cor.thin.posterior} gives
\[
\widehat p_{i,t+1}
=
\frac{
\E{
p_{i,t+1}\left(\Theta_i^{(2)}\right)
\prod_{k=1}^{n_{i\bullet t}}
\Phi\left(u_k\left(h_k+\Theta_i^{(2)}\right)\right)
}
}{
\Phi_{n_{i\bullet t}}
\left(
\mu_{it};
\boldsymbol I_{n_{i\bullet t}}+\sigma^2\boldsymbol U_{it}
\right)
}.
\]
Since
\(p_{i,t+1}(\theta)=\Phi\left(1\cdot(\eta_{i,t+1}+\theta)\right)\), the
product in the numerator is the product in \eqref{eq.thin.D} for the
history augmented by \((n,m)=(1,1)\). Applying Lemma~\ref{lem.bp} to the
augmented history proves part iii.

To show part iv, we first observe that, conditional on
\(\mathcal H_{it}\), the joint density of
\(\left(\Theta_i^{(1)},\Theta_i^{(2)}\right)\) and the probability mass
function of \(N_{i,t+1}\) satisfy
\[
\pi(\theta_1,\theta_2,n\mid\mathcal H_{it})
=
\pi(\theta_1\mid\mathcal H_{it})
\P{N_{i,t+1}=n\mid\Theta_i^{(1)}=\theta_1}
\,
\pi(\theta_2\mid\mathcal H_{it}),
\]
by part i of Corollary~\ref{cor.thin.posterior} and
Model~\ref{model.thin}. Hence \(\Theta_i^{(2)}\) is independent of
\(N_{i,t+1}\) conditional on \(\mathcal H_{it}\). Together with the
binomial observation equation, this gives
\[
\E{M_{i,t+1}\mid\mathcal H_{it},N_{i,t+1}=n}
=
n\,\E{p_{i,t+1}\left(\Theta_i^{(2)}\right)\mid\mathcal H_{it},N_{i,t+1}=n}
=
n\,\widehat p_{i,t+1}.
\]

Finally, part v follows from part iv and part ii, since
\[
\E{M_{i,t+1}\mid\mathcal H_{it}}
=
\E{N_{i,t+1}\mid\mathcal H_{it}}\,\widehat p_{i,t+1}.
\]
\end{proof}

\begin{remark}[Interpretation of the predictions]
\label{rem.thin.interpretation}
Part ii of Theorem~\ref{thm.thin.pred} is the exact credibility formula
of the Poisson-gamma model: the a priori frequency \(\lambda_{i,t+1}\)
is multiplied by a credibility-weighted average of the observed ratio
\(n_{i\bullet t}/\lambda_{i\bullet t}\) and the prior mean \(1\), with
credibility factor \(\zeta_{it}\), which increases to one as
\(\lambda_{i\bullet t}\to\infty\). Part iii is the posterior predictive
probability that an additional Bernoulli-probit trial, namely a new
claim, is a success.
In particular, if \(n_{i\bullet t}=0\), then
\(
\widehat p_{i,t+1}
=
\Phi\left(\eta_{i,t+1}/\sqrt{1+\sigma^2}\right),
\)
which is the a priori probability that a claim is extreme. Part v shows
that the prediction of the extreme claim count is the product of the
two exact Bayesian predictions, one for the frequency and one for the
proportion of extreme claims, as a consequence of the a posteriori
independence of the random effects.
\end{remark}

We conclude this section with the unconditional moments of
Model~\ref{model.thin}, which describe the dependence induced by the
random effects.

\begin{proposition}[Moments]
\label{prop.thin.moments}
Under Model~\ref{model.thin}, fix a region \(i\), and define
\[
\pi_{it}:=\Phi\left(\frac{\eta_{it}}{\sqrt{1+\sigma^2}}\right),
\qquad
\pi_{i,st}:=
\Phi_2\left(
(\eta_{is},\eta_{it})^\top;
\boldsymbol I_2+\sigma^2\boldsymbol J_2
\right),
\qquad s,t\in\mathbb N,
\]
where \(\boldsymbol J_2\) is the \(2\times2\) matrix of ones. Then, for
\(s,t\in\mathbb N\) with \(s\neq t\), we have
\[
\begin{aligned}
\E{N_{it}}&=\lambda_{it},
&
\E{M_{it}}&=\lambda_{it}\pi_{it},\\
\Var{N_{it}}&=\lambda_{it}+\frac{\lambda_{it}^2}{r},
&
\cov{N_{is},N_{it}}&=\frac{\lambda_{is}\lambda_{it}}{r},\\
\cov{N_{it},M_{it}}&=\pi_{it}\left(\lambda_{it}+\frac{\lambda_{it}^2}{r}\right),
&
\cov{N_{is},M_{it}}&=\frac{\lambda_{is}\lambda_{it}\pi_{it}}{r},
\end{aligned}
\]
and
\[
\begin{aligned}
\Var{M_{it}}
&=
\lambda_{it}\pi_{it}
+
\lambda_{it}^2
\left\{
\left(1+\frac1r\right)\pi_{i,tt}-\pi_{it}^2
\right\},\\
\cov{M_{is},M_{it}}
&=
\lambda_{is}\lambda_{it}
\left\{
\left(1+\frac1r\right)\pi_{i,st}-\pi_{is}\pi_{it}
\right\}.
\end{aligned}
\]
\end{proposition}

\begin{proof}
Write \(\Theta_1:=\Theta_i^{(1)}\) and \(\Theta_2:=\Theta_i^{(2)}\), and
recall \(\E{\Theta_1}=1\) and \(\E{\Theta_1^2}=1+1/r\). By
Lemma~\ref{lem.bp} with \(k=1\) and \(k=2\), and \(u_{1:k}\) the
vector of ones,
\[
\E{p_{it}(\Theta_2)}=\pi_{it},
\qquad
\E{p_{is}(\Theta_2)p_{it}(\Theta_2)}=\pi_{i,st},
\qquad s,t\in\mathbb N.
\]
Conditional on \((\Theta_1,\Theta_2)\), Model~\ref{model.thin} and
Lemma~\ref{lem.thin} give, for \(s\neq t\),
\[
\begin{aligned}
\E{N_{it}\mid\Theta_1,\Theta_2}&=\lambda_{it}\Theta_1,
&
\E{M_{it}\mid\Theta_1,\Theta_2}&=\lambda_{it}\Theta_1p_{it}(\Theta_2),\\
\E{N_{it}^2\mid\Theta_1,\Theta_2}&=\lambda_{it}\Theta_1+\lambda_{it}^2\Theta_1^2,
&
\E{M_{it}^2\mid\Theta_1,\Theta_2}&=\lambda_{it}\Theta_1p_{it}(\Theta_2)
+\lambda_{it}^2\Theta_1^2p_{it}(\Theta_2)^2,
\end{aligned}
\]
and
\[
\begin{aligned}
\E{N_{it}M_{it}\mid\Theta_1,\Theta_2}
&=
p_{it}(\Theta_2)\E{N_{it}^2\mid\Theta_1,\Theta_2}
=
p_{it}(\Theta_2)\left(\lambda_{it}\Theta_1+\lambda_{it}^2\Theta_1^2\right),\\
\E{N_{is}M_{it}\mid\Theta_1,\Theta_2}
&=
\lambda_{is}\lambda_{it}\Theta_1^2p_{it}(\Theta_2),\\
\E{M_{is}M_{it}\mid\Theta_1,\Theta_2}
&=
\lambda_{is}\lambda_{it}\Theta_1^2p_{is}(\Theta_2)p_{it}(\Theta_2).
\end{aligned}
\]
Taking expectations, using the independence of \(\Theta_1\) and
\(\Theta_2\), and subtracting the products of the means proves the
result.
\end{proof}
